% ============================================================
%  main.tex — Documento maestro
%  Plantilla corporativa Terabyte para informes y entregables
%  escritos (línea gráfica "técnica contemporánea").
%
%  Compilador: pdfLaTeX  (Overleaf: Menu > Compiler > pdfLaTeX)
%  Estructura: este archivo NO contiene contenido; solo ensambla
%  configuración (config/) + portada (portada/) + secciones
%  (secciones/). El contenido de cada entrega se edita en
%  secciones/, nunca aquí.
% ============================================================
\documentclass[11pt]{article} % cuerpo mínimo 11pt (Bases Administrativas)

% ---------- Configuración base ----------
\input{config/00-preambulo}
\input{config/01-comandos}
\input{portada/portada}

% ============================================================
%  METADATOS DE ESTA ENTREGA — editar en cada nuevo documento
% ============================================================
\establecerPortada
  {Subdocumento 13 · Informe 1}                          % kicker superior
  {Innovaciones}                                         % título principal
  {Terminal Portuario Aconcagua S.A.}                    % subtítulo / cliente
  {Versión 1.0}                                          % versión
  {7 de septiembre de 2026}                              % fecha
  {Propuesta Técnica --- Caso 06 Portuaria. Documento preparado por Terabyte.} % nota

\establecertitulo{Caso 06 --- Portuaria · Informe 1 · Subdoc. 13} % texto pequeño del encabezado (der.)

\hypersetup{
  pdftitle={Subdocumento 13 --- Innovaciones --- Terminal Portuario Aconcagua S.A.},
  pdfauthor={Terabyte},
  pdfsubject={Caso 06 --- Portuaria, Informe 1, Subdocumento 13}
}

% ============================================================
\begin{document}

\portada

% ---------- Páginas preliminares (numeración romana) ----------
\pagenumbering{roman}

\section*{Control de versiones}
\addcontentsline{toc}{section}{Control de versiones}
\markboth{\MakeUppercase{Control de versiones}}{}

\begin{controlversiones}
0.1 & 05-09-2026 & V. Guzmán / M. Reyes & Cartera preliminar de cinco innovaciones y contraste contra el catálogo de requerimientos funcionales. \\
0.2 & 05-09-2026 & V. Guzmán / M. Reyes & Ajuste de alcance al contenido exigido para el Informe 1; verificación de citas normativas contra Bases y Caso. \\
0.3 & 06-09-2026 & V. Guzmán / M. Reyes & Cierre de la innovación de tipo 5; incorporación del criterio de selección de la cartera. \\
1.0 & 07-09-2026 & V. Guzmán / M. Reyes & Versión de entrega del Subdocumento 13 para el Informe 1. \\
\end{controlversiones}

\vspace{6pt}
\noindent\small\textcolor{terabyteGris}{%
Cada entrega debe agregar una fila nueva a esta tabla. No se sobrescriben
versiones anteriores.}

\clearpage

\tableofcontents
\clearpage
% Este subdocumento no contiene figuras; \listoffigures queda desactivado
% para no imprimir un índice vacío. Reactivar si se agregan diagramas.
%\listoffigures
\listoftables
\clearpage

\section*{Glosario y abreviaturas}
\addcontentsline{toc}{section}{Glosario y abreviaturas}
\markboth{\MakeUppercase{Glosario y abreviaturas}}{}

\begin{description}[style=nextline, leftmargin=1.6cm]
  \item[BA] Bases Administrativas del proceso TFEP-01/2026.
  \item[BTT] Bases Técnicas Transversales TFEP-01/2026.
  \item[CP] Bases Técnicas del Caso (Caso 06 --- Portuaria).
  \item[RT] Código de requisito de las Bases Técnicas Transversales
    (por ejemplo, \texttt{RT-26.01}).
  \item[RF] Código de requerimiento funcional del catálogo del
    Subdocumento 3 (por ejemplo, \texttt{RF-REF-08}).
  \item[Alcance 1, 2 y 3] Categorías de emisiones de gases de efecto
    invernadero: directas de la organización, indirectas por energía
    adquirida e indirectas de la cadena de valor, respectivamente.
  \item[CO\textsubscript{2}e] Dióxido de carbono equivalente. Unidad que
    expresa el efecto de distintos gases de efecto invernadero en una
    escala común.
  \item[EDIFACT] Estándar internacional de mensajería electrónica para el
    intercambio de datos entre navieras y terminales.
  \item[GLEC Framework] Marco del \emph{Global Logistics Emissions Council}
    para la contabilización y el reporte de emisiones logísticas.
  \item[Intensidad de carbono] Emisión atribuida por unidad de actividad.
    En este subdocumento, kilogramos de CO\textsubscript{2}e por contenedor
    movilizado.
  \item[Marcha blanca] Período de operación supervisada con datos y usuarios
    reales, en convivencia con la operación vigente y con plan de reversión
    activo, previo al paso a producción.
  \item[Ralentí] Tiempo en que un equipo mantiene el motor encendido sin
    ejecutar trabajo productivo.
  \item[\emph{Reefer}] Contenedor refrigerado. Por extensión, el patio y las
    tomas de conexión que lo alimentan.
  \item[Remoción] Movimiento de patio que reubica un contenedor sin
    trasladarlo a su destino, ejecutado para acceder a otro que está debajo.
  \item[Sello de tiempo] Constancia criptográfica de que un dato existía en
    un instante determinado, verificable por un tercero.
  \item[Simulación de eventos discretos] Técnica de modelado que representa
    la operación como una secuencia de eventos en el tiempo, para evaluar
    escenarios sin intervenir el sistema real.
  \item[Toma] Punto de conexión eléctrica del patio refrigerado al que se
    enchufa un contenedor.
  \item[TRL] \emph{Technology Readiness Level}, nivel de madurez tecnológica.
    En este documento se declara según la escala de la norma ISO 16290:2013.
  \item[Ventana de atraque] Franja horaria comprometida con la naviera para
    la atención de una nave.
\end{description}

\vspace{4pt}
\noindent\small\textcolor{terabyteGris}{%
Agregar cada sigla o término técnico nuevo apenas se use por primera vez
en el cuerpo del documento.}

\clearpage

\section*{Resumen ejecutivo}
\addcontentsline{toc}{section}{Resumen ejecutivo}
\markboth{\MakeUppercase{Resumen ejecutivo}}{}

\kicker{Subdocumento 13 · Innovaciones}

Terabyte presenta la cartera obligatoria de cinco innovaciones exigida por el
Capítulo 5 de las Bases Administrativas, una por cada tipo del Artículo
28\textdegree{}.

La cartera no se compuso eligiendo tecnologías. Se compuso \textbf{descartando}: se
contrastó cada candidata contra los 139 requerimientos funcionales y las 21
decisiones del alcance ya definido, y el resultado fue que la mayor parte de
las respuestas evidentes del sector —sistema de citas, reconocimiento óptico
del contenedor, asignación algorítmica de posición, alarmas de frío
escaladas, mensajería estándar, portal de autoservicio, cálculo de emisiones—
\textbf{ya son alcance obligatorio de esta propuesta}. Presentarlas como
innovación habría activado la causal de descarte del Artículo 30\textdegree{}.

Lo que quedó después de ese filtro, y de las restricciones no negociables del
Capítulo 10 del Caso, no son tecnologías sino \textbf{cinco huecos} entre lo
que el terminal necesita como resultado y lo que los requerimientos cubren
como conducta:

\begin{enumerate}[leftmargin=*, itemsep=2pt]
  \item El sistema avisa hacia adentro y nunca hacia afuera: cuando una toma
    refrigerada falla, el dueño de la carga no se entera.
  \item La cita de camión existe, pero el proceso sigue suponiendo que el
    transportista controla cuándo estará lista la carga.
  \item No existe dónde ensayar un cambio: cuatro meses y medio de
    congelamiento, 620 recaladas al año y un patio al 90\,\% en temporada.
  \item El CLIENTE compra el equipamiento de terreno, no tiene quién lo
    mantenga y su proveedor no comparte la consecuencia del indicador
    comprometido con el concedente.
  \item Los requerimientos obligan a \emph{medir} las emisiones. Ninguno
    obliga a \emph{reducirlas}.
\end{enumerate}

Cada innovación se desarrolla en su idea, la tecnología que la sustenta, su
alcance, su forma de implementación y su resultado esperado, conforme al
Formulario T-22, y declara expresamente aquello que requiere investigación
adicional. Al menos una de ellas es verificable durante la marcha blanca de la
Etapa 1, antes del mes 16, conforme al \texttt{RT-26.08}.

\begin{bloquePropuesta}
Ninguna de las cinco es transferible a otro caso del llamado: cada una está
anclada a un hecho medido de este terminal. Y ninguna es indispensable para
cumplir el alcance comprometido: todas degradan a la funcionalidad obligatoria
si no rinden lo esperado.
\end{bloquePropuesta}

\noindent\small\textcolor{terabyteGris}{%
Conforme al Artículo 53\textdegree{} de las Bases Administrativas, este
subdocumento no contiene información económica. La valorización de las cinco
innovaciones en el flujo de caja corresponde al Informe 3.}

\clearpage

% ---------- Cuerpo del informe (numeración arábiga) ----------
\pagenumbering{arabic}
\setcounter{page}{1}

% Cada subdocumento se entrega como PDF independiente (comunicado del
% profesor, punto 6). En este proyecto solo se compila el Subdocumento 13.
%\input{secciones/01-subdoc}\clearpage
%\input{secciones/02-subdoc}\clearpage
%\input{secciones/03-subdoc}\clearpage
%\input{secciones/04-subdoc}\clearpage
%\input{secciones/05-subdoc}\clearpage

% El Subdocumento 13 debe numerarse "13" aunque sea la única sección:
\setcounter{section}{12}
% ============================================================
%  secciones/13-subdoc.tex
%  SUBDOCUMENTO 13 — INNOVACIONES
%  Célula 4 (V. Guzmán / M. Reyes) — Informe 1, 7 de septiembre de 2026
%
%  NOTA SOBRE CITAS
%  La plantilla ofrece \citaCP{p}, \citaBA{p} y \citaBTT{p}, que exigen
%  el NÚMERO DE PÁGINA del PDF oficial. Este archivo cita por capítulo,
%  artículo y código RT, que es la referencia estable y verificable que
%  usa el equipo. Cuando se disponga de los PDF paginados, reemplazar
%  cada cita textual por el comando correspondiente. NO inventar páginas.
%
%  NUMERACIÓN
%  Esta sección debe numerarse como "13". En main.tex, antes del \input,
%  debe ir:   \setcounter{section}{12}
% ============================================================

\section{Innovaciones}

\kicker{Subdocumento 13}

\noindent
Terabyte presenta la cartera obligatoria de cinco innovaciones, una por cada
tipo del Artículo 28\textdegree{} de las Bases Administrativas. Conforme al
Formulario T-22, en este Informe 1 cada innovación se desarrolla en su idea,
la tecnología que la sustenta, su alcance, su forma de implementación y su
resultado esperado, y declara expresamente aquello que requiere investigación
adicional.

\subsection{Criterio de selección y coherencia con el caso}

\subsubsection{Cómo se eligieron estas cinco y no otras}

El Artículo 30.2 advierte que \enquote{una innovación de alta sofisticación
técnica que no resuelva un problema real del caso obtendrá menor puntaje que
una innovación sencilla, bien justificada y con impacto verificable}. Terabyte
tomó esa advertencia como método de trabajo y no como recomendación. La cartera
se construyó mediante un embudo de cuatro descartes sucesivos, aplicado en este
orden.

\paragraph*{Primer descarte: lo que las Bases prohíben expresamente.}
El Artículo 30\textdegree{} y el Capítulo 26 de las Bases Técnicas
Transversales coinciden en la misma exclusión: no se acepta como innovación una
tecnología que ya constituye estándar de la industria, una tendencia sin diseño
de incorporación, ni una funcionalidad exigida por las Bases Técnicas
presentada como innovación. Este último criterio fue el más restrictivo.
Terabyte contrastó cada candidata contra los 139 requerimientos funcionales y
las 21 decisiones del alcance definido en el Subdocumento 3, uno por uno. El
resultado fue que la mayor parte de las respuestas evidentes —sistema de citas,
reconocimiento óptico del contenedor en el gate, asignación algorítmica de
posición, alarmas de frío con escalamiento, mensajería EDIFACT, portal de
autoservicio, cálculo de emisiones por contenedor, credencial temporal de
eventuales— \textbf{ya son alcance obligatorio}, y presentarlas como innovación
habría sido entregar al evaluador la causal de descarte.

\paragraph*{Segundo descarte: las restricciones no negociables.}
La Restricción N\textdegree{}~1 —ninguna solución puede aumentar la exposición
de una persona al riesgo del patio— y la Restricción N\textdegree{}~11 —el área
de tecnologías de información del CLIENTE son cinco personas— eliminan buena
parte del catálogo tecnológico habitual de la industria: automatización de
equipos de patio, vehículos guiados, realidad aumentada en cabina, inventario
por vehículo aéreo. El gerente general dejó consignada en el acta del
directorio la condición de que \enquote{cualquier solución que aumente la
exposición al riesgo en el patio será rechazada, por buena que sea en todo lo
demás} \Hecho.

\paragraph*{Tercer descarte: las exclusiones explícitas del Capítulo 11.}
Queda fuera toda innovación que suponga reemplazar el sistema de gestión
empresarial, intervenir el control de las grúas de muelle, sustituir la
videovigilancia, operar la báscula, electrificar la flota o ejecutar obra
civil.

\paragraph*{Cuarto filtro: anclaje a un hecho medido de este terminal.}
Cada innovación que sobrevivió debe apuntar a una cifra del Capítulo 7 o a un
vacío declarado del numeral 16.1. Una innovación que no pueda hacerlo no es
pertinente a este caso, por buena que sea en abstracto.

\subsubsection{Qué quedó después del embudo}

Lo que sobrevive no son tecnologías: son cinco huecos entre lo que el caso
exige como resultado y lo que los requerimientos cubren como conducta. Cada uno
cae, además, en un tipo distinto del Artículo 28\textdegree{}, lo que permite
componer la cartera obligatoria sin forzar ninguna pieza.

\begin{table}[htbp]
\centering
\caption{Huecos identificados y su correspondencia con los tipos obligatorios}
\label{tab:huecos}
\begin{tabularx}{\textwidth}{@{}>{\raggedright\arraybackslash}p{2.9cm}>{\raggedright\arraybackslash}X>{\raggedright\arraybackslash}p{1.3cm}>{\raggedright\arraybackslash}p{2.9cm}@{}}
\toprule
\rowcolor{terabyteFondo}
\textcolor{white}{\bfseries Hueco} &
\textcolor{white}{\bfseries Enunciado} &
\textcolor{white}{\bfseries ID} &
\textcolor{white}{\bfseries Tipo Art. 28\textdegree{}} \\
\midrule
De destinatario &
El sistema avisa hacia adentro y nunca hacia afuera. La alarma de frío llega al
operador y al supervisor; nadie avisa al dueño de la carga &
IN-01 & 1 · Producto o servicio \\
\addlinespace
De proceso &
La cita existe, pero el proceso sigue suponiendo que el transportista controla
cuándo estará lista la carga &
IN-02 & 2 · Proceso \\
\addlinespace
De condición del terreno &
No existe dónde ensayar un cambio. No es un indicador: es una restricción
estructural, y ningún requerimiento la atiende &
IN-03 & 3 · Tecnología o arquitectura \\
\addlinespace
Contractual &
El CLIENTE compra el equipamiento, no tiene quién lo mantenga y su proveedor no
comparte la consecuencia del indicador comprometido &
IN-04 & 4 · Modelo de negocio \\
\addlinespace
De verbo &
Los requerimientos dicen \emph{medir}. Ninguno dice \emph{reducir} &
IN-05 & 5 · Sostenibilidad \\
\bottomrule
\end{tabularx}
\end{table}

\begin{bloquePropuesta}
La cartera se construyó desde el problema hacia la tecnología, y no al revés.
Ninguna de las cinco innovaciones de Terabyte es transferible a otro caso del
llamado: cada una está anclada a un hecho de este terminal —el evento del 18 de
febrero, el transportista que no controla cuándo estará lista su carga, el
congelamiento estacional de cuatro meses y medio, las cinco personas de TI y la
fecha de 2029 con el 34\,\% de los contenedores.
\end{bloquePropuesta}

\subsubsection{Diferencia entre innovación y funcionalidad obligatoria}

El Artículo 30\textdegree{} de las Bases Administrativas y el Capítulo 26 de
las Bases Técnicas Transversales prohíben lo mismo con la misma frase:
\enquote{No se aceptará como innovación: la sola adopción de una tecnología que
ya constituye estándar de la industria; la mención de una tendencia sin diseño
de incorporación; ni una funcionalidad exigida por las Bases Técnicas
presentada como innovación.}

Terabyte aplicó ese filtro de forma verificable y declara aquí el resultado del
contraste, requerimiento por requerimiento, para que la Comisión pueda
comprobarlo.

\begin{longtable}{@{}>{\raggedright\arraybackslash}p{1.5cm}>{\raggedright\arraybackslash}p{3.1cm}>{\raggedright\arraybackslash}p{4.6cm}>{\raggedright\arraybackslash}p{6.2cm}@{}}
\caption{Contraste de cada innovación contra el requerimiento obligatorio más
próximo}\label{tab:nosolapamiento}\\
\toprule
\rowcolor{terabyteFondo}
\textcolor{white}{\bfseries ID} &
\textcolor{white}{\bfseries Requerimiento más próximo} &
\textcolor{white}{\bfseries Qué obliga ese requerimiento} &
\textcolor{white}{\bfseries Qué agrega la innovación} \\
\midrule
\endfirsthead
\toprule
\rowcolor{terabyteFondo}
\textcolor{white}{\bfseries ID} &
\textcolor{white}{\bfseries Requerimiento más próximo} &
\textcolor{white}{\bfseries Qué obliga ese requerimiento} &
\textcolor{white}{\bfseries Qué agrega la innovación} \\
\midrule
\endhead
\bottomrule
\endfoot
IN-01 & \texttt{RF-REF-08}, \texttt{RF-REF-11}, \texttt{RF-REF-12} &
Notificar la alarma al operador de turno y a un supervisor de guardia; mantener
la serie continua; permitir descargar la serie por autoservicio &
El destinatario del aviso pasa a ser el dueño de la carga, que hoy no recibe
nada; y la serie deja de ser un archivo descargable para convertirse en un
certificado sellado, verificable por un tercero sin depender del terminal \\
\addlinespace
IN-02 & \texttt{RF-GAT-01}, \texttt{RF-GAT-02}, \texttt{RF-GAT-16}, Decisión
N\textdegree{}~6 &
Solicitar cita; priorizar a quien la cumple; limitar franjas por capacidad
declarada; cola virtual que avisa al transportista cuando el terminal se atrasa &
Incorpora al dueño de la carga como tercera parte de la cita y reprograma la
franja automáticamente cuando confirma disponibilidad. Resuelve la mitad de la
decisión pendiente N\textdegree{}~6 del numeral 16.1 que ninguna decisión
resolvió \\
\addlinespace
IN-03 & \texttt{RF-NAV-05}, \texttt{RF-GAT-16} &
Estimar la duración de la nave a partir de serie histórica; limitar franjas por
capacidad declarada &
Ningún requerimiento dice de dónde sale la capacidad declarada ni permite
ensayar un cambio de política antes de aplicarlo en producción. No existe
requerimiento de simulación en el catálogo \\
\addlinespace
IN-04 & Ninguno & --- &
El catálogo no contiene ningún requerimiento sobre esquema contractual ni sobre
modelo de servicio. IN-04 no compite con ningún requerimiento funcional \\
\addlinespace
IN-05 & \texttt{RF-EMI-01} a \texttt{RF-EMI-06} &
Medir el consumo, calcular la emisión por contenedor, trazar el cálculo,
acumular la serie y obtener el reporte verificado &
Ningún requerimiento obliga a \emph{reducir}. IN-05 compromete una meta de
intensidad y actúa sobre remociones evitadas, ralentí y asignación preferente
de los dos equipos eléctricos existentes \\
\end{longtable}

\subsubsection{Coherencia conjunta de la cartera}

Las cinco no son iniciativas sueltas: comparten infraestructura y se sostienen
unas a otras.

\begin{itemize}[leftmargin=*, itemsep=3pt]
  \item \textbf{IN-03 alimenta a IN-02.} El gemelo produce la capacidad
    declarada de franjas que exige \texttt{RF-GAT-16} y permite verificar,
    antes de tocar el gate real, si la reprogramación a tres bandas mejora o
    empeora la fila.
  \item \textbf{IN-01 e IN-05 comparten la cadena de evidencia.} Ambas dependen
    del servicio de evidencia y firma y de la trazabilidad hasta el dato de
    origen; una la usa para temperatura, la otra para emisiones. El sellado se
    construye una vez.
  \item \textbf{IN-04 depende de que IN-02 e IN-05 sean medibles.} Un
    compromiso de resultado solo es defendible si el indicador que lo verifica
    lo produce el propio sistema de forma trazable y reproducible por un
    tercero, que es lo que exige \texttt{RF-CON-11} y la Restricción
    N\textdegree{}~14.
  \item \textbf{Ninguna toca lo excluido.} Ninguna de las cinco interviene el
    control de grúas de muelle, la emisión tributaria del sistema de gestión
    empresarial, la videovigilancia ni la báscula.
  \item \textbf{Ninguna añade interacción del operador en el patio.} La
    Restricción N\textdegree{}~1 se respeta por diseño, no por declaración.
\end{itemize}

\subsubsection{Relación con las etapas y con el alcance}

El reparto sigue el alcance definido en el Subdocumento 3 y el cronograma
contractual obligatorio de 56 meses del Artículo 17\textdegree{} \Hecho.

\begin{table}[htbp]
\centering
\caption{Materialización de la cartera en el cronograma contractual}
\label{tab:etapas}
\begin{tabularx}{\textwidth}{@{}>{\raggedright\arraybackslash}p{4.3cm}>{\raggedright\arraybackslash}p{1.6cm}>{\raggedright\arraybackslash}X@{}}
\toprule
\rowcolor{terabyteFondo}
\textcolor{white}{\bfseries Fase contractual} &
\textcolor{white}{\bfseries Meses} &
\textcolor{white}{\bfseries Qué de la cartera ocurre allí} \\
\midrule
Etapa 1 · Desarrollo & 1 a 12 &
IN-05 inicia captura de consumo desde el mes 1; el cálculo de intensidad existe
aquí solo como prototipo y prevalidación. IN-03 v1 se construye entre los meses
6 y 12. IN-02 se construye sobre el gate de Etapa 1 \\
\addlinespace
Etapa 1 · Marcha blanca & 13 a 15 &
IN-02 se mide aquí, con lo que la propuesta satisface el \texttt{RT-26.08}
deseable. IN-01 fase A opera con aviso por canal directo \\
\addlinespace
Etapa 1 · Producción & 16 &
IN-04 entra en vigencia con la primera medición oficial \\
\addlinespace
Etapa 2 · Desarrollo & 13 a 18 &
IN-01 fase B, certificado autenticado por portal. IN-03 v2 \\
\addlinespace
Etapa 2 · Marcha blanca & 19 a 20 & IN-01 fase B en convivencia \\
\addlinespace
Etapa 2 · Producción & 21 &
IN-05 entra con motor de cálculo productivo y activa el término de reducción.
IN-04 en régimen pleno \\
\addlinespace
Operación & 21 a 56 &
Las cinco en régimen; IN-05 acumula la trayectoria que se presenta a la alianza
naviera en 2029 \\
\bottomrule
\end{tabularx}
\end{table}

\begin{bloqueSupuesto}
La asignación de cada innovación a paquetes de la estructura de descomposición
del trabajo y las fechas finas del cronograma se establecen en el Subdocumento
7 y en el Formulario T-14, y se presentan en el Informe 2, conforme a la
progresión que fija el Formulario T-22. Los meses indicados en el cuadro
anterior son la hipótesis de trabajo anclada al Artículo 17\textdegree{}.
\end{bloqueSupuesto}

% ============================================================
\subsection{Innovación de producto o servicio: cadena de frío certificada
Aconcagua}
% ============================================================

\paragraph*{Idea.}
Hoy, cuando una toma refrigerada falla, el terminal se avisa a sí mismo. El
exportador cuya fruta está en esa toma no se entera. Y cuando la carga llega a
destino, el terminal no puede entregar el registro continuo de temperatura que
los mercados exigen como evidencia de cadena de frío: de las 2.400 tomas de
conexión del patio refrigerado y de sus 26 tableros, ninguno tiene hoy
instrumentación remota, y el control se hace mediante una ronda a pie con
planilla cada cuatro horas \Hecho. La innovación convierte la instrumentación que el proyecto
instalará de todos modos en un servicio con dos caras nuevas: el aviso llega al
dueño de la carga en el momento en que ocurre la desviación, y la serie de
temperatura se entrega como un certificado sellado que un tercero puede
verificar por su cuenta, sin pedirle nada al terminal \Propuesta.

\paragraph*{Tecnología que la sustenta.}
Tres piezas, ninguna experimental. La cadena de eventos de temperatura por
contenedor se sella mediante registro encadenado por función resumen, se firma
electrónicamente y recibe sello de tiempo conforme a RFC~3161, de modo que la
fecha y la integridad son demostrables sin acceso a los sistemas del terminal.
La pertenencia de un registro individual a la serie sellada se demuestra
mediante estructura de árbol de Merkle, siguiendo el patrón de registro
auditable de RFC~6962, sin exponer el resto de la serie. El aviso externo viaja
por los adaptadores de canal por audiencia que la solución ya contempla, con
confirmación de recepción registrada. La firma reutiliza el servicio
transversal que la solución ya construye para los cuatro actos obligatorios: no
se agrega ninguna plataforma criptográfica nueva.

\paragraph*{Alcance.}
Cubre los contenedores refrigerados conectados a tomas instrumentadas. Se
despliega en dos fases: primero el aviso proactivo por canal directo, y después
el certificado descargable con usuario autenticado, porque el registro y la
verificación de identidad de clientes externos pertenecen a la Etapa 2. No
cubre la carga seca, ni la temperatura durante el transporte marítimo, ni
sustituye el registro del propio equipo refrigerado del contenedor.

\paragraph*{Forma de implementación.}
Se apoya íntegramente en componentes que el alcance obligatorio ya construye:
el contexto de reefer y telemetría genera el evento; el servicio de evidencia y
firma lo sella; el servicio de notificaciones lo emite por el adaptador que
corresponda a cada audiencia; el portal y la puerta de enlace de servicios
entregan el certificado y exponen el punto público de verificación. Al abrir
una nueva superficie de exposición externa, requiere modelado de amenazas
propio conforme al \texttt{RT-26.07}, que se desarrolla en el Subdocumento 4.

\paragraph*{Resultado esperado.}
Que una desviación de temperatura llegue a alguien que tenga interés propio en
que se resuelva, en cualquiera de los tres turnos: exactamente lo que faltó la
madrugada del 18 de febrero de 2026, cuando la pérdida fue de US\$~620.000 en
38 contenedores durante nueve horas \Hecho. Y que el terminal pueda ofrecer
evidencia de cadena de frío como atributo de su servicio, en una operación
donde entre enero y marzo se concentra el 62\,\% del volumen refrigerado del
año. El Caso registra la consecuencia comercial de no tenerla en cuatro
palabras: tras aquella pérdida, \enquote{el exportador no volvió}.

\begin{bloqueDuda}
\textbf{Investigación adicional declarada.} Requiere levantar, con al menos un
exportador y una autoridad, qué formato de certificado es efectivamente
aceptado en los mercados de destino de la fruta chilena. Sin esa validación, el
certificado puede ser técnicamente correcto y comercialmente inútil. La meta
del indicador se fijará una vez conocido ese formato.
\end{bloqueDuda}

% ============================================================
\subsection{Innovación de proceso: cita convenida a tres bandas}
% ============================================================

\paragraph*{Idea.}
Un sistema de citas para camiones ordena las llegadas, pero no cambia el hecho
que las desordena. Lo dijo con todas sus letras el gerente general de una
empresa de transporte terrestre cuyo 30\,\% de operación pasa por este
terminal, en la entrevista de levantamiento:

\begin{quote}
\small
\enquote{Yo no controlo cuándo está lista la carga. El packing me llama y me
dice \enquote{ya, ven}. Si me dan una cita a las diez y la fruta está a las
tres, la cita no sirve y voy a llegar igual a las tres. [\dots] si van a hacer
citas, tienen que ser citas que se puedan cambiar y que consideren que la mitad
de mi operación no depende de mí.}
\end{quote}

Pedirle una hora al transportista es pedirle una promesa que no depende de él,
y el Caso dejó esa pregunta explícitamente sin resolver en su lista de
decisiones pendientes. La innovación cambia el proceso de negocio: la cita deja
de ser una solicitud del transportista y pasa a ser un acuerdo de tres partes
—dueño de la carga, transportista y terminal— en el que la franja se puede
cambiar sola, porque se reprograma automáticamente cuando el dueño de la carga
confirma que la carga está efectivamente disponible \Propuesta. Es, literalmente,
la cita que el transportista pidió.

\paragraph*{Tecnología que la sustenta.}
No es una tecnología nueva, y el Artículo 28\textdegree{} admite expresamente
que una innovación de proceso lo sea. El modelo tiene cuatro movimientos: la
reserva se asocia a una unidad de carga y no solo a una patente; el dueño de la
carga confirma disponibilidad por el portal o por integración, y esa
confirmación libera la franja; si la confirmación no llega dentro del plazo
declarado, el sistema reprograma automáticamente a la siguiente franja
compatible y notifica a las tres partes; la franja liberada se reasigna, de
modo que la capacidad declarada del gate no se pierde.

\paragraph*{Alcance.}
Aplica a los movimientos de retiro y recepción coordinados con un dueño de
carga identificable. No aplica a movimientos sin contraparte externa
identificable, que siguen el flujo de cita simple. No introduce penalización, y
esa decisión también viene de la entrevista: \enquote{Si me penalizan por no
cumplir una cita que nunca pude cumplir, en dos semanas nadie va a usar el
sistema.}

\paragraph*{Forma de implementación.}
El contexto de gate y citas incorpora el estado \enquote{franja condicionada a
confirmación} y la regla de reprogramación; el portal expone la confirmación al
dueño de la carga; el servicio de notificaciones avisa a las tres partes por
los adaptadores ya previstos. Se suma un rol externo a la matriz de
autorización del servicio de identidad. No modifica la arquitectura de
seguridad.

\paragraph*{Resultado esperado.}
Atacar la estadía del camión —78 minutos contra los 45 comprometidos con el
concedente, con tres semestres consecutivos sobre el umbral, y una fila que el
12 de marzo de 2026 alcanzó 3,2 kilómetros y 140 camiones \Hecho— por una vía
distinta de la que atacan los requerimientos obligatorios: en vez de acelerar
la atención dentro del terminal, evitar que el camión salga a la ruta antes de
que su carga esté lista. Es la innovación con la que Terabyte satisface el
\texttt{RT-26.08}, que valora que al menos una sea verificable durante la
marcha blanca de la Etapa 1, antes del mes 16.

\begin{bloqueDuda}
\textbf{Investigación adicional declarada.} La variante de tres partes con
reprogramación automática no está documentada en la literatura arbitrada
revisada; lo que sí está documentado es el mecanismo base de los sistemas de
cita. Se declara como adaptación de una práctica madura a una condición no
resuelta de este caso, y no como tecnología probada. Requiere además levantar
con transportistas y exportadores el plazo realista de confirmación.
\end{bloqueDuda}

% ============================================================
\subsection{Innovación tecnológica o de arquitectura: gemelo de operación del
Terminal Aconcagua}
% ============================================================

\paragraph*{Idea.}
En este terminal no existe dónde ensayar un cambio. Está prohibido intervenir
entre el 15 de diciembre y el 30 de abril; está prohibido intervenir durante la
atención de una nave y en las cuatro horas previas a una ventana confirmada,
con 620 recaladas al año; el patio opera al 90\,\% en temporada y la operación
es 24x7x365 sin ventana de detención \Hecho. En consecuencia, toda decisión
operacional relevante se toma hoy sin poder probarla, y se prueba en producción
con una nave amarrada. La innovación es un banco de decisiones fuera de línea:
un modelo del terminal, calibrado con la telemetría que el propio proyecto
instala, donde el cambio se ensaya antes de tocar la operación \Propuesta.

\paragraph*{Tecnología que la sustenta.}
Simulación de eventos discretos del ciclo completo —llegada de camión, gate,
asignación de posición, movimiento de patio, tractocamión, grúa de muelle,
ventana de nave— calibrada contra la telemetría real de equipos y los eventos
de gate que la solución produce de todos modos. La calibración es lo que lo
distingue de un modelo de escritorio: el gemelo se contrasta contra la
operación observada y declara su error junto a cada resultado. En su versión
comprometida no incorpora inteligencia artificial; si se incorporara, quedaría
sujeta íntegramente al Capítulo 18 de las Bases Técnicas Transversales,
conforme al \texttt{RT-26.06}.

\paragraph*{Alcance.}
Versión 1: gate y flujo del camión. Versión 2: patio, nave y planificación.
Una restricción de alcance es parte de la innovación: \textbf{el gemelo no
tiene autoridad operacional}. No expone ninguna interfaz de escritura hacia los
contextos operacionales, no emite órdenes y su indisponibilidad no degrada
ninguna función. Se despliega en nube y nunca en el borde, precisamente para no
competir por la capacidad local que sostiene las cinco funciones críticas
durante las 72 horas sin enlace.

\paragraph*{Forma de implementación.}
Componente nuevo, alimentado exclusivamente desde la capa analítica y desde las
series temporales de telemetría. Se incorpora al catálogo de componentes del
Subdocumento 4, con su emplazamiento justificado conforme al Artículo
16\textdegree{}. La interdicción de escritura se verifica en la matriz de
segregación de funciones del servicio de identidad, de modo que sea auditable y
no una declaración.

\paragraph*{Resultado esperado.}
Tres cosas concretas. Que cada cambio de política de patio o de gate que hoy
solo puede probarse con una nave amarrada pase a probarse en el modelo. Que la
capacidad declarada que limita las franjas de cita —un parámetro que los
requerimientos exigen pero cuyo origen ninguno define— tenga por fin un
fundamento. Y que la pregunta del Caso sobre qué componente se satura primero
en el peak y ante la eventual incorporación de un cuarto sitio entre 2030 y
2032 se responda con un modelo en vez de con una opinión.

\begin{bloqueDuda}
\textbf{Investigación adicional declarada.} El nivel de madurez se declara como
rango —TRL 6 a 7 en la escala ISO 16290:2013— y no como valor único, porque
depende de la calidad de la calibración, que a su vez depende de
instrumentación que hoy no existe: solo 12 de los 18 equipos de patio tienen
terminal montada. La meta de error del gemelo no puede fijarse antes de conocer
la densidad real de la telemetría instalada.
\end{bloqueDuda}

% ============================================================
\subsection{Innovación de modelo de negocio o contratación: servicio gestionado
del equipamiento de terreno}
% ============================================================

\paragraph*{Idea.}
El caso plantea una asimetría que ningún requerimiento funcional puede
resolver, porque no es técnica sino contractual. El CLIENTE compra el hardware
pero no tiene quién lo mantenga: el proponente especifica qué comprar, el
CLIENTE lo adquiere, y su área de tecnologías de información son cinco personas
en un operador de importancia vital \Hecho. Ese equipamiento va además a un
ambiente que lo destruye —atmósfera salina, corrosión acelerada, a menos de
trescientos metros de la línea de costa—, y las Bases exigen que su plan de
reposición esté costeado. La innovación es que Terabyte gestione el ciclo de
vida completo de ese equipamiento como servicio, y que asuma un compromiso de
resultado verificable sobre indicadores que hoy nadie garantiza \Propuesta.

\paragraph*{Tecnología que la sustenta.}
Contratación basada en desempeño, aplicada con dos componentes acoplados.
Primero, un servicio gestionado de ciclo de vida del equipamiento de terreno:
monitoreo de estado, mantención preventiva, reposición por corrosión e
inventario de repuestos, sobre equipamiento que permanece en el patrimonio del
CLIENTE. Segundo, un compromiso de resultado sobre indicadores que el propio
sistema produce de forma trazable, con dos salvaguardas de diseño: ningún
indicador comprometido puede calcularse con un dato que Terabyte produzca sin
evidencia reproducible por un tercero, y el compromiso recae sobre lo que
Terabyte controla.

\paragraph*{Alcance, y su límite frente a las Bases.}
El compromiso se refiere a indicadores que el Artículo 78\textdegree{} no
cubre: disponibilidad del equipamiento de terreno instrumentado —tomas,
tableros y equipos móviles— y continuidad de la instrumentación en ambiente
marino. El Artículo 78\textdegree{} regula disponibilidad del servicio, tiempos
de respuesta y resolución, tiempo medio de restauración, tasa de cambios
fallidos, cumplimiento de los objetivos de punto y tiempo de recuperación,
vulnerabilidades, reincidencia y satisfacción; no regula el equipamiento de
terreno del CLIENTE.

\begin{bloquePropuesta}
\textbf{Declaración expresa de Terabyte.} Esta innovación \textbf{se suma} al
régimen de niveles de servicio del Artículo 78\textdegree{}, a la medición del
Artículo 79\textdegree{} y al régimen de multas del Artículo 80\textdegree{}, y
\textbf{en ningún caso los reemplaza, los rebaja ni propone un esquema
alternativo de pago, de multas o de topes}. Terabyte acepta íntegramente y sin
reserva alguna el régimen contractual de las Bases.
\end{bloquePropuesta}

\paragraph*{Forma de implementación.}
No es un componente de software. Su viabilidad descansa en la cadena de
evidencia que la solución ya construye: el cálculo trazable de la estadía, la
evidencia de alarma y confirmación de frío, la emisión trazable hasta el dato
de origen y la producción trazable de los indicadores del concedente. La
medición se realiza sobre la plataforma de observabilidad, cuyos datos el
Artículo 79\textdegree{} ya obliga a poner a disposición del CLIENTE en tiempo
real y de forma exportable.

\paragraph*{Resultado esperado.}
Que la Restricción N\textdegree{}~11 —\enquote{toda función que requiera un
especialista dedicado que la compañía no tiene debe ofrecerse como servicio y
estar costeada}— se cumpla en su letra y no en abstracto; que el plan de
reposición del equipamiento en atmósfera marina que la Restricción
N\textdegree{}~12 exige sea exigible a alguien; y que el proveedor comparta la
consecuencia de un indicador que el terminal lleva tres semestres consecutivos
incumpliendo ante su concedente.

\begin{bloqueDuda}
\textbf{Investigación adicional declarada.} Falta definir el conjunto exacto de
indicadores comprometidos y el momento de activación, y falta resolver el
problema de atribución: la estadía del camión depende de factores que Terabyte
no controla —volumen comercial, ocupación del patio, comportamiento del
transportista—, de modo que el compromiso debe recaer sobre indicadores de
disponibilidad y de comportamiento del sistema, con un mecanismo declarado de
exclusión por causas ajenas documentadas.
\end{bloqueDuda}

% ============================================================
\subsection{Innovación de sostenibilidad: intensidad de carbono comprometida}
% ============================================================

\paragraph*{Idea.}
La alianza naviera impuso tres condiciones para 2029, y una es el reporte
verificado de emisiones \Hecho. El alcance obligatorio resuelve la mitad del
problema: medir, calcular por contenedor, trazar hasta el dato de origen,
acumular la serie y obtener el reporte verificado. Pero medir no es reducir. Lo
que la alianza recibiría en 2029, si el proyecto solo cumple lo obligatorio, es
una cifra sin trayectoria. La innovación es comprometer una meta de intensidad
de carbono —kilogramos de CO\textsubscript{2} equivalente por contenedor
movilizado— y reducirla con la misma metodología que el verificador va a
auditar \Propuesta.

\paragraph*{Tecnología que la sustenta.}
La misma norma del reporte obligatorio, usada como instrumento de gestión y no
solo de reporte: ISO 14083:2023 como norma de cuantificación, implementada
operativamente mediante el GLEC Framework v3.2 y verificada por tercero
acreditado bajo ISO 14064-3. Sobre esa base, tres palancas concretas:
\emph{remociones evitadas} —cada remoción es un movimiento que quema
combustible sin trasladar carga, y hoy son el 18\,\% de los movimientos de
patio—; \emph{ralentí de tractocamiones}, derivado de la telemetría de consumo
que la solución instala de todos modos en los 16 equipos diésel, sin sensor
adicional; y \emph{asignación preferente de los dos equipos eléctricos
existentes} a los ciclos de mayor intensidad de emisión.

\paragraph*{Alcance.}
Cubre las emisiones de alcance 1 y 2 atribuibles a los movimientos de
contenedores dentro del terminal, con el alcance 3 explicitado. No electrifica
la flota, respetando la exclusión expresa del Caso, y no compromete emisión
absoluta sino intensidad por contenedor, precisamente para que la meta no
dependa del volumen comercial, que el terminal no controla.

\paragraph*{Forma de implementación.}
El contexto de energía y emisiones incorpora el motor de intensidad y la
trayectoria contra meta; el contexto de patio incorpora un término de
puntuación de emisiones dentro del algoritmo de asignación que ya existe, con
peso configurable y tope declarado. Dos restricciones de diseño son parte de la
innovación: el término de emisiones nunca desplaza una restricción de
seguridad, una regla de segregación de mercancías peligrosas ni una ventana de
nave comprometida; y no agrega ninguna interacción del operador, con lo que la
Restricción N\textdegree{}~1 se respeta por construcción.

\paragraph*{Resultado esperado.}
Que el terminal llegue a 2029 ante el 34\,\% de sus contenedores no con una
cifra sino con una trayectoria de reducción verificada bajo la misma norma; que
cada exportador y cada naviera reciban el dato de emisión atribuido a su
embarque, utilizable en su propio inventario de alcance 3; y que las remociones
evitadas y el ralentí eliminado se traduzcan en combustible no comprado.
Responde además a la objeción que la gerenta comercial dejó registrada:
\enquote{si empezamos la mensajería y las emisiones el 2028, no llegamos,
porque un reporte de emisiones verificado necesita serie de datos previa.}

\begin{bloqueDuda}
\textbf{Investigación adicional declarada.} La captura de consumo comienza en
el mes 1, pero el motor de cálculo productivo entra con la producción de la
Etapa 2, en el mes 21; lo que exista antes es prototipo y prevalidación
metodológica con el verificador, y así se declara. \textbf{La meta de reducción
no puede fijarse hoy y no se finge que sí:} la línea base del Caso es literal,
las emisiones por contenedor \enquote{no se miden}, y no existe porcentaje de
reducción comprometible contra una línea base inexistente. Terabyte compromete,
en su lugar, declarar la línea base como entregable verificable y fijar la meta
en ese momento, con acuerdo del verificador acreditado. Falta además someter
tempranamente al verificador el método de atribución del consumo al contenedor
individual, que es el punto donde un verificador puede objetar.
\end{bloqueDuda}

% ============================================================
\subsection{Trazabilidad conjunta de las innovaciones}
% ============================================================

\subsubsection{Resultados verificables esperados}

\begin{longtable}{@{}>{\raggedright\arraybackslash}p{1.3cm}>{\raggedright\arraybackslash}p{4.4cm}>{\raggedright\arraybackslash}p{3.6cm}>{\raggedright\arraybackslash}p{6.1cm}@{}}
\caption{Indicador de verificación y línea base por
innovación}\label{tab:indicadores}\\
\toprule
\rowcolor{terabyteFondo}
\textcolor{white}{\bfseries ID} &
\textcolor{white}{\bfseries Indicador de verificación} &
\textcolor{white}{\bfseries Línea base} &
\textcolor{white}{\bfseries Momento de medición} \\
\midrule
\endfirsthead
\toprule
\rowcolor{terabyteFondo}
\textcolor{white}{\bfseries ID} &
\textcolor{white}{\bfseries Indicador de verificación} &
\textcolor{white}{\bfseries Línea base} &
\textcolor{white}{\bfseries Momento de medición} \\
\midrule
\endhead
\bottomrule
\endfoot
IN-01 &
Embarques refrigerados con certificado emitido y verificable; tiempo entre la
desviación y el aviso recibido por el dueño de la carga &
Cero: hoy el dueño de la carga no recibe aviso alguno y no existe registro
continuo entregable &
Marcha blanca de Etapa 1 (fase A) y de Etapa 2 (fase B) \\
\addlinespace
IN-02 &
Citas cumplidas dentro de ventana sobre citas otorgadas; reprogramaciones
originadas por confirmación del dueño de la carga; estadía con cita convenida &
Cero: no existe sistema de citas. Estadía general: 78 minutos &
Semanal durante la marcha blanca de Etapa 1, meses 13 a 15 \\
\addlinespace
IN-03 &
Error del gemelo contra la operación real en estadía y movimientos por hora;
cambios de política ensayados antes de aplicarse &
Cero: no existe capacidad de ensayo ni modelo de comparación &
Calibración al cierre de desarrollo de Etapa 1; error medido en cada marcha
blanca \\
\addlinespace
IN-04 &
Indicadores comprometidos alcanzados en el período; disponibilidad del
equipamiento de terreno gestionado &
71\,\% de cumplimiento de ventana de atraque; tres semestres de estadía sobre
umbral &
Semestral, alineado al informe de indicadores al concedente \\
\addlinespace
IN-05 &
Intensidad de carbono en kg CO\textsubscript{2}e por contenedor movilizado,
conforme a ISO 14083:2023 &
\textbf{No se mide} &
Prototipo hacia el mes 12; motor productivo en el mes 21; verificación bajo
ISO 14064-3 antes de 2029 \\
\end{longtable}

\begin{bloqueSupuesto}
\textbf{Por qué cuatro metas quedan declaradas y no fijadas.} El Capítulo 18
del Caso obliga al proponente a \enquote{proponer la meta cuando este documento
no la fije}. Terabyte aplicó ese criterio con rigor en el Subdocumento 3,
documentando meta por meta la evidencia que la sostiene e incluso la que le es
adversa. Comprometer aquí cuatro metas sin ese mismo trabajo produciría números
que no resistirían la defensa técnica. Conforme al Formulario T-22, esta
propuesta declara la investigación pendiente antes que ofrecer una cifra que no
puede sostener.
\end{bloqueSupuesto}

\subsubsection{Riesgos de adopción de la cartera}

Además del riesgo propio de cada innovación, la cartera tiene tres riesgos
comunes que conviene declarar juntos.

\begin{riesgo}[R-INN-01]
\enuncioriesgo{IN-01 e IN-02 requieren que un actor externo —el dueño de la
carga— haga algo que hoy no hace, y su adopción no depende del
ADJUDICATARIO}{la adopción se estanque en niveles bajos}{el beneficio
diferenciador no se materialice, aunque el alcance obligatorio se cumpla
igualmente}
\end{riesgo}

\begin{riesgo}[R-INN-02]
\enuncioriesgo{cuatro de las cinco innovaciones se materializan en meses que el
congelamiento estacional puede desplazar, y está prohibido intervenir entre el
15 de diciembre y el 30 de abril, durante la atención de una nave y en las
cuatro horas previas a una ventana confirmada}{una ventana de despliegue se
pierda}{la verificación del beneficio se desplace a la temporada siguiente}
\end{riesgo}

\begin{riesgo}[R-INN-03]
\enuncioriesgo{el Artículo 30.3 establece que las innovaciones comprometidas
forman parte del alcance contractual y son exigibles como cualquier otro
requerimiento}{una innovación no rinda lo esperado}{su omisión sea tratada como
incumplimiento del alcance}
\end{riesgo}

\begin{bloquePropuesta}
\textbf{Mitigación transversal.} Cada innovación de esta cartera declara un
plan de contingencia que degrada a la funcionalidad obligatoria sin dejar un
compromiso incumplido. \textbf{Ningún requerimiento obligatorio depende de
ninguna de las cinco innovaciones}: si una no rinde, el alcance comprometido se
cumple igualmente y solo se pierde el valor adicional.
\end{bloquePropuesta}

\subsubsection{Alcance de esta entrega}

Conforme al Formulario T-22, en este Informe 1 las cinco innovaciones se
presentan nombradas y desarrolladas en su idea, la tecnología que las sustenta,
su alcance, su forma de implementación y su resultado esperado, y cada una
declara expresamente aquello que requiere investigación adicional.

El Formulario T-19 con los siete elementos del Artículo 29\textdegree{} se
presenta con la propuesta final. El refinamiento del alcance trazable con la
arquitectura y con la estructura de descomposición del trabajo corresponde al
Informe 2, y la valorización de cada innovación en el flujo de caja
—inversión, costo operacional y beneficio esperado— corresponde al Informe 3.
Por esa razón, y conforme al Artículo 53\textdegree{} de las Bases
Administrativas, \textbf{este subdocumento no contiene información económica}.
\clearpage

% Para reproducir saltos de numeración oficial (p. ej. el Informe 1
% exige que la sección de innovaciones se numere "13" aunque sea la
% sexta sección del documento), forzar el contador antes de \section:
%
%   \setcounter{section}{12}
%   \input{secciones/05-innovaciones}

% Este PDF es el Subdocumento 13 entregado por separado (comunicado del
% profesor, punto 6). Las "Conclusiones del Informe 1" pertenecen al informe
% consolidado, no a un subdocumento individual: quedan desactivadas.
%\input{secciones/general/97-conclusiones}
\section*{Referencias bibliográficas}
\addcontentsline{toc}{section}{Referencias bibliográficas}
\markboth{\MakeUppercase{Referencias bibliográficas}}{}

\noindent\small\textcolor{terabyteGris}{%
Formato APA 7.\textsuperscript{a} edición, conforme al Artículo
29\textdegree{} de las Bases Administrativas y al \texttt{RT-26.03} de las
Bases Técnicas Transversales, que exigen citar las fuentes de las
innovaciones de base tecnológica.}

\vspace{10pt}
\noindent\textbf{\small Fuentes de las innovaciones}

\vspace{4pt}
\begin{list}{}{\leftmargin=1.2cm\itemindent=-1.2cm\itemsep=5pt}
  \item Adams, C., Cain, P., Pinkas, D., \& Zuccherato, R. (2001).
    \emph{Internet X.509 public key infrastructure time-stamp protocol (TSP)}
    (RFC 3161). Internet Engineering Task Force.
    \url{https://doi.org/10.17487/RFC3161}

  \item Carlo, H. J., Vis, I. F. A., \& Roodbergen, K. J. (2014). Storage yard
    operations in container terminals: Literature overview, trends, and
    research directions. \emph{European Journal of Operational Research,
    235}(2), 412--430.

  \item Giuliano, G., \& O'Brien, T. (2007). Reducing port-related truck
    emissions: The terminal gate appointment system at the Ports of Los
    Angeles and Long Beach. \emph{Transportation Research Part D: Transport
    and Environment}.

  \item Hakimi, F., Khaled, T., Al-Kharaz, M., Foahom Gouabou, A. C., \&
    Amzil, K. (2024). Towards a digital twin modeling method for container
    terminal port. \emph{Procedia Computer Science, 246}, 3113--3121.
    \url{https://doi.org/10.1016/j.procs.2024.09.361}

  \item International Organization for Standardization. (2013). \emph{Space
    systems --- Definition of the technology readiness levels (TRLs) and their
    criteria of assessment} (ISO 16290:2013).

  \item International Organization for Standardization. (2019).
    \emph{Greenhouse gases --- Part 3: Specification with guidance for the
    verification and validation of greenhouse gas statements}
    (ISO 14064-3:2019).

  \item International Organization for Standardization. (2023).
    \emph{Greenhouse gases --- Quantification and reporting of greenhouse gas
    emissions arising from transport chain operations} (ISO 14083:2023).

  \item Laurie, B., Langley, A., \& Kasper, E. (2013). \emph{Certificate
    transparency} (RFC 6962). Internet Engineering Task Force.
    \url{https://doi.org/10.17487/RFC6962}

  \item Ramírez-Nafarrate, A., González-Ramírez, R. G., Smith, N. R.,
    Guerra-Olivares, R., \& Voß, S. (2017). Impact on yard efficiency of a
    truck appointment system for a port terminal. \emph{Annals of Operations
    Research, 258}(2), 195--216.
    \url{https://doi.org/10.1007/s10479-016-2384-0}

  \item Selviaridis, K., \& Wynstra, F. (2015). Performance-based contracting:
    A literature review and future research directions. \emph{International
    Journal of Production Research, 53}(12), 3505--3540.
    \url{https://doi.org/10.1080/00207543.2014.978031}

  \item Smart Freight Centre. (2023). \emph{Global Logistics Emissions Council
    framework for logistics emissions accounting and reporting} (Version 3.0).
\end{list}

\vspace{10pt}
\noindent\textbf{\small Fuentes normativas del proceso}

\vspace{4pt}
\begin{list}{}{\leftmargin=1.2cm\itemindent=-1.2cm\itemsep=5pt}
  \item Pontificia Universidad Católica de Valparaíso. (2026).
    \emph{FEP01.26 Bases Administrativas TFEP-01-2026}. Artículos
    16\textdegree{}, 17\textdegree{}, 28\textdegree{}, 29\textdegree{},
    30\textdegree{}, 53\textdegree{}, 78\textdegree{}, 79\textdegree{} y
    80\textdegree{}; Formularios T-19, T-21 y T-22.
  \item Pontificia Universidad Católica de Valparaíso. (2026).
    \emph{FEP02.26 Bases Técnicas Transversales TFEP-01-2026}. Capítulos 18,
    22 y 26 (\texttt{RT-26.01} a \texttt{RT-26.08}).
  \item Pontificia Universidad Católica de Valparaíso. (2026).
    \emph{FEP03.06.26 Caso 06 --- Portuaria (Bases Técnicas del Caso)}.
    Capítulos 1, 6, 7, 8, 9, 10, 11, 12, 13, 14, 16, 17, 18 y 19.
\end{list}

\vspace{8pt}
\begin{bloqueDuda}
La referencia de Giuliano y O'Brien (2007) está confirmada en autoría, título,
año y revista, pero su volumen, número y rango de páginas deben completarse
contra el original antes de la entrega final. No se consignan datos no
verificados.
\end{bloqueDuda}


% ---------- Anexos ----------
% Los anexos A--K pertenecen al informe consolidado, no a este subdocumento.
% Además, el Formulario T-19 completo se exige en la propuesta final y no en
% el Informe 1, por lo que aquí no corresponde anexarlo. Bloque desactivado.
%\clearpage
%\appendix
%\enanexotrue
%\renewcommand{\thesection}{\Alph{section}}
%\setcounter{section}{0}
%\input{secciones/general/99-anexos}

\end{document}
